\chapter{Substructures}

\section{Vector sublattices}

A \emph{vector sublattice} of a vector lattice $X$ is a linear subspace
closed under the lattice operations. Since the lattice operations are
inter-definable, closure under any one of them is enough: closure under
$\vee$ or under $\lvert \cdot \rvert$ both characterise the notion. This
section introduces the bundled structure, shows that every vector
sublattice is a vector lattice in its own right and that the collection
of vector sublattices of a fixed ambient space forms a complete lattice
under inclusion, and describes the vector sublattice generated by a set
in several concrete situations — a submodule, a pointed cone, and a
pairwise-disjoint family of non-negative vectors. The section closes
with the normed setting: every vector sublattice of a normed vector
lattice is itself a normed vector lattice; its norm closure is again a
vector sublattice; norm-closed vector sublattices of a Banach lattice
are themselves Banach lattices; and the closed vector sublattice
generated by a countable subset is separable.

\subsection{Pointed cones}

Before introducing vector sublattices themselves, we record that in a
vector lattice the sup-closure and inf-closure of a pointed cone are
again pointed cones. These two operations will appear as basic building
blocks in the descriptions of the generated vector sublattice. Recall
that a \emph{pointed cone} is a subset of a vector subspace that is
closed under addition, positive scalar multiplication, and contains
the origin.

\begin{lemma}
\label{lem:pointedCone-supClosure}
\lean{PointedCone.supClosure}
\leanok
\uses{def:vector-lattice}
Let $X$ be a vector lattice and let $C$ be a pointed cone in $X$. Then
the sup-closure of $C$, viewed as a subset of $X$, is a pointed cone.
\end{lemma}

\begin{lemma}
\label{lem:pointedCone-infClosure}
\lean{PointedCone.infClosure}
\leanok
\uses{def:vector-lattice}
Let $X$ be a vector lattice and let $C$ be a pointed cone in $X$. Then
the inf-closure of $C$, viewed as a subset of $X$, is a pointed cone.
\end{lemma}

\subsection{Vector sublattices}

Throughout the rest of the section $X$ denotes a vector lattice.

\begin{definition}
\label{def:vector-sublattice}
\lean{VectorSublattice}
\leanok
\uses{def:vector-lattice}
Let $X$ be a vector lattice. A \emph{vector sublattice} of $X$ is a
linear subspace $Y \subseteq X$ which is closed under $\vee$: for all
$x, y \in Y$, $x \vee y \in Y$.
\end{definition}

\begin{lemma}
\label{lem:vectorSublattice-inf-mem}
\lean{VectorSublattice.inf_mem}
\leanok
\uses{def:vector-sublattice, lem:inf-eq-sub-posPart}
Let $Y$ be a vector sublattice of $X$ and let $x, y \in Y$. Then
$x \wedge y \in Y$.
\end{lemma}

\begin{lemma}
\label{lem:vectorSublattice-posPart-mem}
\lean{VectorSublattice.posPart_mem}
\leanok
\uses{def:vector-sublattice}
Let $Y$ be a vector sublattice of $X$ and let $x \in Y$. Then
$x^+ \in Y$.
\end{lemma}

\begin{lemma}
\label{lem:vectorSublattice-negPart-mem}
\lean{VectorSublattice.negPart_mem}
\leanok
\uses{lem:vectorSublattice-posPart-mem}
Let $Y$ be a vector sublattice of $X$ and let $x \in Y$. Then
$x^- \in Y$.
\end{lemma}

\begin{lemma}
\label{lem:vectorSublattice-abs-mem}
\lean{VectorSublattice.abs_mem}
\leanok
\uses{def:vector-sublattice}
Let $Y$ be a vector sublattice of $X$ and let $x \in Y$. Then
$\lvert x \rvert \in Y$.
\end{lemma}

\begin{lemma}
\label{lem:vectorSublattice-isSublattice}
\lean{VectorSublattice.isSublattice}
\leanok
\uses{def:vector-sublattice, lem:vectorSublattice-inf-mem}
Every vector sublattice of $X$ is a sublattice of $X$.
\end{lemma}

Since the lattice operations of a vector lattice are inter-definable with
the absolute value, a linear subspace closed under $\lvert \cdot \rvert$
is automatically a vector sublattice.

\begin{lemma}
\label{lem:vectorSublattice-ofAbsClosed}
\lean{VectorSublattice.ofAbsClosed}
\leanok
\uses{def:vector-sublattice, lem:sup-eq-add-posPart}
A linear subspace of $X$ closed under $\lvert \cdot \rvert$ (i.e.\
$\lvert x \rvert \in M$ for every $x \in M$) is a vector sublattice of
$X$.
\end{lemma}

\begin{lemma}
\label{lem:vectorSublattice-abs-mem-iff-sup-mem}
\lean{VectorSublattice.abs_mem_iff_sup_mem}
\leanok
\uses{lem:vectorSublattice-ofAbsClosed, lem:vectorSublattice-abs-mem}
Let $M$ be a linear subspace of $X$. Then $M$ is closed under
$\lvert \cdot \rvert$ if and only if $M$ is closed under $\vee$.
\end{lemma}

A vector sublattice inherits a vector lattice structure from the ambient
space, with the lattice operations, addition and scalar multiplication
computed pointwise.

\begin{lemma}
\label{lem:vectorSublattice-instVectorLatticeSubtype}
\lean{VectorSublattice.instVectorLatticeSubtype}
\leanok
\uses{def:vector-lattice, def:vector-sublattice, lem:vectorSublattice-inf-mem}
Let $Y$ be a vector sublattice of $X$. Then $Y$ is itself a vector
lattice, with lattice operations, addition, and scalar multiplication
inherited pointwise from $X$.
\end{lemma}

The collection of vector sublattices of $X$ is ordered by inclusion. It
has a greatest element $\top = X$, a least element $\bot = \{0\}$, and is
closed under arbitrary intersections.

\begin{lemma}
\label{lem:vectorSublattice-toSubmodule-injective}
\lean{VectorSublattice.toSubmodule_injective}
\leanok
\uses{def:vector-sublattice}
The forgetful map from vector sublattices of $X$ to linear subspaces of
$X$ is injective.
\end{lemma}

\begin{lemma}
\label{lem:vectorSublattice-complete-lattice}
\lean{VectorSublattice.instTop, VectorSublattice.instBot, VectorSublattice.instInfSet}
\leanok
\uses{def:vector-sublattice, lem:vectorSublattice-toSubmodule-injective}
The whole space $X$ is the greatest vector sublattice of $X$, the zero
subspace $\{0\}$ is the least vector sublattice of $X$, and the
intersection of an arbitrary family of vector sublattices is again a
vector sublattice of $X$.
\end{lemma}

\subsection{The vector sublattice generated by a set}

For any subset $s \subseteq X$, there is a smallest vector sublattice of
$X$ containing $s$, obtained as the infimum of all vector sublattices
that contain $s$. The subsection characterises this generated sublattice
in several concrete situations.

\begin{definition}
\label{def:vectorSublattice-generated}
\lean{VectorSublattice.generated}
\leanok
\uses{lem:vectorSublattice-complete-lattice}
Let $s$ be a subset of $X$. The vector sublattice \emph{generated} by
$s$ is the intersection of all vector sublattices of $X$ containing $s$.
\end{definition}

\begin{lemma}
\label{lem:vectorSublattice-subset-generated}
\lean{VectorSublattice.subset_generated}
\leanok
\uses{def:vectorSublattice-generated}
For every subset $s \subseteq X$, $s$ is contained in the vector
sublattice generated by $s$.
\end{lemma}

\begin{lemma}
\label{lem:vectorSublattice-generated-le}
\lean{VectorSublattice.generated_le}
\leanok
\uses{def:vectorSublattice-generated}
Let $s \subseteq X$ and let $Y$ be a vector sublattice of $X$ with
$s \subseteq Y$. Then the vector sublattice generated by $s$ is
contained in $Y$.
\end{lemma}

\begin{theorem}
\label{thm:vectorSublattice-generated-submodule-eq-supClosure-infClosure}
\lean{VectorSublattice.generated_submodule_eq_supClosure_infClosure}
\leanok
\uses{def:vectorSublattice-generated, lem:nonneg-smul-sup, lem:nonneg-smul-inf}
Let $M$ be a linear subspace of $X$. Then the vector sublattice
generated by $M$ equals the sup-closure of the inf-closure of $M$:
\[
  \operatorname{generated}(M) = \operatorname{supClosure}(
  \operatorname{infClosure}(M)).
\]
\end{theorem}

\begin{theorem}
\label{thm:vectorSublattice-generated-submodule-eq-infClosure-supClosure}
\lean{VectorSublattice.generated_submodule_eq_infClosure_supClosure}
\leanok
\uses{thm:vectorSublattice-generated-submodule-eq-supClosure-infClosure}
Let $M$ be a linear subspace of $X$. Then the vector sublattice
generated by $M$ equals the inf-closure of the sup-closure of $M$:
\[
  \operatorname{generated}(M) = \operatorname{infClosure}(
  \operatorname{supClosure}(M)).
\]
\end{theorem}

\begin{theorem}
\label{thm:vectorSublattice-generated-pointedCone-eq-sub-supClosure}
\lean{VectorSublattice.generated_pointedCone_eq_sub_supClosure}
\leanok
\uses{def:vectorSublattice-generated, lem:pointedCone-supClosure, lem:vectorSublattice-isSublattice}
Let $C$ be a pointed cone in $X$. Then the vector sublattice generated
by $C$ equals the difference set of the sup-closure of $C$ with
itself:
\[
  \operatorname{generated}(C) = \operatorname{supClosure}(C)
  - \operatorname{supClosure}(C).
\]
\end{theorem}

\begin{theorem}
\label{thm:vectorSublattice-generated-pointedCone-eq-sub-infClosure}
\lean{VectorSublattice.generated_pointedCone_eq_sub_infClosure}
\leanok
\uses{def:vectorSublattice-generated, lem:pointedCone-infClosure, lem:vectorSublattice-isSublattice}
Let $C$ be a pointed cone in $X$. Then the vector sublattice generated
by $C$ equals the difference set of the inf-closure of $C$ with
itself:
\[
  \operatorname{generated}(C) = \operatorname{infClosure}(C)
  - \operatorname{infClosure}(C).
\]
\end{theorem}

\begin{theorem}
\label{thm:vectorSublattice-generated-eq-span-of-pairwise-vlDisjoint}
\lean{VectorSublattice.generated_eq_span_of_pairwise_vlDisjoint}
\leanok
\uses{def:vectorSublattice-generated, lem:vectorSublattice-ofAbsClosed, thm:abs-sum-of-pairwise-isVLDisjoint}
Let $A \subseteq X$ be a set of pairwise-disjoint elements with
$0 \le a$ for every $a \in A$. Then the vector sublattice generated by
$A$ coincides with the linear span of $A$ over $\mathbb{R}$.
\end{theorem}

\subsection{Normed and Banach sublattices}

In the normed setting, every vector sublattice of a normed vector
lattice is itself a normed vector lattice under the inherited norm and
lattice operations. The topological closure of a vector sublattice is
again a vector sublattice; in a Banach lattice, norm-closed vector
sublattices are themselves Banach lattices. Finally, the closed
vector sublattice generated by a countable subset is separable.

\begin{lemma}
\label{lem:vectorSublattice-instNormedVectorLatticeSubtype}
\lean{VectorSublattice.instNormedVectorLatticeSubtype}
\leanok
\uses{def:normed-vector-lattice, def:vector-sublattice, lem:vectorSublattice-instVectorLatticeSubtype}
Let $X$ be a normed vector lattice and let $Y$ be a vector sublattice of
$X$. Then $Y$ is itself a normed vector lattice under the induced norm
and lattice operations.
\end{lemma}

\begin{lemma}
\label{lem:vectorSublattice-topologicalClosure}
\lean{VectorSublattice.topologicalClosure}
\leanok
\uses{def:vector-sublattice, def:normed-vector-lattice, thm:continuous-sup}
Let $X$ be a normed vector lattice and let $Y$ be a vector sublattice of
$X$. Then the norm closure of $Y$ is a vector sublattice of $X$.
\end{lemma}

\begin{lemma}
\label{lem:vectorSublattice-instBanachLatticeSubtype}
\lean{VectorSublattice.instBanachLatticeSubtype}
\leanok
\uses{def:banach-lattice, def:vector-sublattice, lem:vectorSublattice-instNormedVectorLatticeSubtype}
Let $X$ be a Banach lattice and let $Y$ be a norm-closed vector
sublattice of $X$. Then $Y$ is a Banach lattice under the induced structures.
\end{lemma}

\begin{theorem}
\label{thm:vectorSublattice-separableSpace-topologicalClosure-generated-of-countable}
\lean{VectorSublattice.separableSpace_topologicalClosure_generated_of_countable}
\leanok
\uses{lem:vectorSublattice-topologicalClosure, def:vectorSublattice-generated, thm:vectorSublattice-generated-submodule-eq-supClosure-infClosure}
Let $X$ be a normed vector lattice and let $A \subseteq X$ be countable.
Then the topological closure of the vector sublattice generated by $A$
is separable.
\end{theorem}

\section{Ideals}

An \emph{order ideal} of a vector lattice $X$ is a vector sublattice
$J$ of $X$ which is \emph{solid}: whenever $x \in J$ and $0 \le y \le
x$, $y$ lies in $J$ as well. Equivalently, a linear subspace $M$ of
$X$ is solid in the absolute-value sense
$\lvert y \rvert \le \lvert x \rvert \Rightarrow y \in M$ if and only
if it is closed under $\vee$ and solid in the positive-cone sense;
in either formulation the resulting notion coincides with that of a
solid subspace. This section develops the basic theory of order
ideals: it introduces the bundled structure, produces the principal
ideal generated by a single vector and equips it, in the Archimedean
setting, with a canonical gauge (or order-unit) norm that makes it a
normed vector lattice, shows that the ideals of $X$ ordered by
inclusion form a complete lattice, produces a proper non-trivial
ideal whenever the dimension of $X$ is strictly greater than $1$, and studies
the topological closure of an ideal in a normed vector lattice
together with the resulting notion of closed order ideal in a
Banach lattice.

\subsection{Order ideals}

Throughout the rest of the section $X$ denotes a vector lattice.

\begin{definition}
\label{def:order-ideal}
\lean{OrderIdeal}
\leanok
\uses{def:vector-sublattice}
Let $X$ be a vector lattice. An \emph{order ideal} of $X$ is a vector
sublattice $J$ of $X$ which is \emph{solid}: for all $x, y \in X$
with $x \in J$, $0 \le y$ and $y \le x$, $y \in J$.
\end{definition}

\begin{lemma}
\label{lem:orderIdeal-mem-of-abs-mem}
\lean{OrderIdeal.mem_of_abs_mem}
\leanok
\uses{def:order-ideal}
Let $J$ be an order ideal of $X$ and let $x \in X$. If
$\lvert x \rvert \in J$, then $x \in J$.
\end{lemma}

\begin{lemma}
\label{lem:orderIdeal-mem-of-abs-le-abs}
\lean{OrderIdeal.mem_of_abs_le_abs}
\leanok
\uses{def:order-ideal, lem:orderIdeal-mem-of-abs-mem}
Let $J$ be an order ideal of $X$ and let $x, y \in X$ with $x \in J$
and $\lvert y \rvert \le \lvert x \rvert$. Then $y \in J$.
\end{lemma}

\begin{lemma}
\label{lem:orderIdeal-ofSolid}
\lean{OrderIdeal.ofSolid}
\leanok
\uses{def:order-ideal}
A linear subspace of $X$ which is solid in the absolute-value sense
(for all $x, y \in X$ with $x \in M$ and $\lvert y \rvert \le
\lvert x \rvert$, $y \in M$) is an order ideal of $X$.
\end{lemma}

\begin{lemma}
\label{lem:orderIdeal-solid-iff}
\lean{OrderIdeal.solid_iff}
\leanok
\uses{lem:orderIdeal-ofSolid}
Let $M$ be a linear subspace of $X$. Then $M$ is solid in the sense
of the absolute value (for all $x, y \in X$ with $x \in M$ and
$\lvert y \rvert \le \lvert x \rvert$, $y \in M$) if and only if $M$
is both solid in the positive-cone sense (for all $x, y \in X$ with
$x \in M$, $0 \le y$ and $y \le x$, $y \in M$) and closed under
$\vee$ (for all $x, y \in M$, $x \vee y \in M$).
\end{lemma}

\begin{lemma}
\label{lem:orderIdeal-toSubmodule-injective}
\lean{OrderIdeal.toSubmodule_injective}
\leanok
\uses{def:order-ideal, lem:vectorSublattice-toSubmodule-injective}
The forgetful map from order ideals of $X$ to linear subspaces of
$X$ is injective.
\end{lemma}

\subsection{The principal ideal}

Every vector $a \in X$ determines an order ideal $I_a$ of $X$, the
\emph{principal ideal} of $a$, consisting of the vectors whose
absolute value is dominated by a non-negative scalar multiple of
$\lvert a \rvert$.

\begin{definition}
\label{def:orderIdeal-principal}
\lean{OrderIdeal.principal}
\leanok
\uses{def:order-ideal}
Let $a \in X$. The \emph{principal ideal} generated by $a$ is the
order ideal
\[
  I_a := \{x \in X : \exists\, c \in \mathbb{R},\ 0 \le c \text{ and }
    \lvert x \rvert \le c\,\lvert a \rvert\}
\]
of $X$.
\end{definition}

\begin{lemma}
\label{lem:orderIdeal-self-mem-principal}
\lean{OrderIdeal.self_mem_principal}
\leanok
\uses{def:orderIdeal-principal}
Let $a \in X$. Then $a \in I_a$.
\end{lemma}

Every vector in the principal ideal $I_a$ carries a canonical
non-negative quantity: the infimum of all non-negative scalars
$c$ for which $\lvert x \rvert \le c\,\lvert a \rvert$. This is the
\emph{gauge norm} (or \emph{order-unit norm}) of $x$ with respect to
$a$. The gauge norm is absolutely homogeneous on $X$, and
subadditive and monotone in $\lvert \cdot \rvert$ on the principal
ideal; together with its definiteness — which requires $X$ to be
Archimedean — this makes $I_a$ a normed vector lattice whenever
$a \ne 0$.

\begin{definition}
\label{def:orderIdeal-gaugeNorm}
\lean{OrderIdeal.gaugeNorm}
\leanok
\uses{def:vector-lattice}
Let $a, x \in X$. The \emph{gauge norm} of $x$ with respect to $a$
is
\[
  \lVert x \rVert_a := \inf\{c \in \mathbb{R} : 0 \le c \text{ and }
    \lvert x \rvert \le c\,\lvert a \rvert\},
\]
with the convention that the infimum of an empty subset of
$\mathbb{R}$ is $0$.
\end{definition}

\begin{lemma}
\label{lem:orderIdeal-gaugeNorm-nonneg}
\lean{OrderIdeal.gaugeNorm_nonneg}
\leanok
\uses{def:orderIdeal-gaugeNorm}
Let $a, x \in X$. Then $0 \le \lVert x \rVert_a$.
\end{lemma}

\begin{lemma}
\label{lem:orderIdeal-gaugeNorm-zero}
\lean{OrderIdeal.gaugeNorm_zero}
\leanok
\uses{def:orderIdeal-gaugeNorm}
Let $a \in X$. Then $\lVert 0 \rVert_a = 0$.
\end{lemma}

\begin{lemma}
\label{lem:orderIdeal-gaugeNorm-neg}
\lean{OrderIdeal.gaugeNorm_neg}
\leanok
\uses{def:orderIdeal-gaugeNorm}
Let $a, x \in X$. Then $\lVert -x \rVert_a = \lVert x \rVert_a$.
\end{lemma}

\begin{lemma}
\label{lem:orderIdeal-gaugeNorm-add-le}
\lean{OrderIdeal.gaugeNorm_add_le}
\leanok
\uses{def:orderIdeal-gaugeNorm, def:orderIdeal-principal}
Let $a \in X$ and let $x, y \in I_a$. Then
$\lVert x + y \rVert_a \le \lVert x \rVert_a + \lVert y \rVert_a$.
\end{lemma}

\begin{lemma}
\label{lem:orderIdeal-gaugeNorm-smul}
\lean{OrderIdeal.gaugeNorm_smul}
\leanok
\uses{def:orderIdeal-gaugeNorm, thm:abs-smul}
Let $a, x \in X$ and let $r \in \mathbb{R}$. Then
$\lVert r \cdot x \rVert_a = \lvert r \rvert \, \lVert x \rVert_a$.
\end{lemma}

\begin{theorem}
\label{thm:orderIdeal-abs-le-gaugeNorm-smul-abs}
\lean{OrderIdeal.abs_le_gaugeNorm_smul_abs}
\leanok
\uses{def:orderIdeal-gaugeNorm, def:orderIdeal-principal,
  def:is-vl-archimedean}
Let $X$ be an Archimedean vector lattice, let $a \in X$ and let
$x \in I_a$. Then
$\lvert x \rvert \le \lVert x \rVert_a \cdot \lvert a \rvert$.
\end{theorem}

\begin{lemma}
\label{lem:orderIdeal-gaugeNorm-le-of-abs-le}
\lean{OrderIdeal.gaugeNorm_le_of_abs_le}
\leanok
\uses{def:orderIdeal-gaugeNorm}
Let $a, x \in X$ and let $c \in \mathbb{R}$ with $0 \le c$ and
$\lvert x \rvert \le c\,\lvert a \rvert$. Then
$\lVert x \rVert_a \le c$.
\end{lemma}

\begin{lemma}
\label{lem:orderIdeal-gaugeNorm-mono-abs}
\lean{OrderIdeal.gaugeNorm_mono_abs}
\leanok
\uses{def:orderIdeal-gaugeNorm, def:orderIdeal-principal}
Let $a \in X$, let $y \in I_a$ and let $x \in X$ with
$\lvert x \rvert \le \lvert y \rvert$. Then
$\lVert x \rVert_a \le \lVert y \rVert_a$.
\end{lemma}

\begin{lemma}
\label{lem:orderIdeal-gaugeNorm-eq-zero-iff}
\lean{OrderIdeal.gaugeNorm_eq_zero_iff}
\leanok
\uses{thm:orderIdeal-abs-le-gaugeNorm-smul-abs,
  lem:abs-eq-zero-iff-zero}
Let $X$ be an Archimedean vector lattice, let $a \in X$ and let
$x \in I_a$. Then $\lVert x \rVert_a = 0$ if and only if $x = 0$.
\end{lemma}

\begin{lemma}
\label{lem:orderIdeal-gaugeNorm-le-one-iff}
\lean{OrderIdeal.gaugeNorm_le_one_iff}
\leanok
\uses{thm:orderIdeal-abs-le-gaugeNorm-smul-abs,
  lem:orderIdeal-gaugeNorm-le-of-abs-le}
Let $X$ be an Archimedean vector lattice, let $a \in X$ and let
$x \in I_a$. Then $\lVert x \rVert_a \le 1$ if and only if
$\lvert x \rvert \le \lvert a \rvert$.
\end{lemma}

\begin{lemma}
\label{lem:orderIdeal-gaugeNorm-anti}
\lean{OrderIdeal.gaugeNorm_anti}
\leanok
\uses{def:orderIdeal-gaugeNorm, def:orderIdeal-principal}
Let $e, u \in X$ with $0 \le e$ and $e \le u$, and let $x \in I_e$.
Then $\lVert x \rVert_u \le \lVert x \rVert_e$.
\end{lemma}

\begin{theorem}
\label{thm:orderIdeal-principalNormedVectorLattice}
\lean{OrderIdeal.principalNormedVectorLattice}
\leanok
\uses{def:orderIdeal-principal, def:orderIdeal-gaugeNorm,
  def:normed-vector-lattice, def:is-vl-archimedean,
  lem:orderIdeal-gaugeNorm-add-le, lem:orderIdeal-gaugeNorm-smul,
  lem:orderIdeal-gaugeNorm-eq-zero-iff,
  lem:orderIdeal-gaugeNorm-mono-abs}
Let $X$ be an Archimedean vector lattice and let $a \in X$ with
$a \ne 0$. Then $I_a$, equipped with the gauge norm
$\lVert \cdot \rVert_a$ and with the lattice and vector-space
operations inherited pointwise from $X$, is a normed vector lattice.
\end{theorem}

\subsection{The ideal generated by a set}

For any subset $s \subseteq X$ there is a smallest order ideal of
$X$ containing $s$, the \emph{ideal generated} by $s$, defined as
the intersection of all order ideals containing $s$. It admits a
concrete description: an element belongs to it precisely when its
absolute value is dominated by a non-negative linear combination of
absolute values of elements of $s$. The principal ideal $I_a$ is
recovered as the ideal generated by the singleton $\{\lvert a
\rvert\}$.

\begin{definition}
\label{def:orderIdeal-generated}
\lean{OrderIdeal.generated}
\leanok
\uses{def:order-ideal}
Let $s \subseteq X$. The order ideal \emph{generated} by $s$ is the
intersection of all order ideals of $X$ containing $s$.
\end{definition}

\begin{lemma}
\label{lem:orderIdeal-subset-generated}
\lean{OrderIdeal.subset_generated}
\leanok
\uses{def:orderIdeal-generated}
Let $s \subseteq X$. Then $s$ is contained in the order ideal
generated by $s$.
\end{lemma}

\begin{lemma}
\label{lem:orderIdeal-generated-le}
\lean{OrderIdeal.generated_le}
\leanok
\uses{def:orderIdeal-generated}
Let $s \subseteq X$ and let $J$ be an order ideal of $X$ with
$s \subseteq J$. Then the ideal generated by $s$ is contained in
$J$.
\end{lemma}

\begin{theorem}
\label{thm:orderIdeal-mem-generated-iff}
\lean{OrderIdeal.mem_generated_iff}
\leanok
\uses{def:orderIdeal-generated, lem:orderIdeal-ofSolid,
  lem:orderIdeal-mem-of-abs-mem}
Let $s \subseteq X$ and let $x \in X$. Then $x$ belongs to the
order ideal generated by $s$ if and only if there exist a finite
set $t \subseteq s$ and non-negative scalars $(c_y)_{y \in t}$
such that
\[
  \lvert x \rvert \le \sum_{y \in t} c_y \, \lvert y \rvert.
\]
\end{theorem}

\begin{lemma}
\label{lem:orderIdeal-principal-eq-generated-abs}
\lean{OrderIdeal.principal_eq_generated_abs}
\leanok
\uses{def:orderIdeal-principal, def:orderIdeal-generated,
  lem:orderIdeal-generated-le, lem:orderIdeal-mem-of-abs-mem}
Let $a \in X$. Then the principal ideal $I_a$ equals the order
ideal generated by the singleton $\{\lvert a \rvert\}$.
\end{lemma}

\subsection{Sum of ideals}

The sum $J_1 + J_2$ of two order ideals, viewed as linear subspaces
of $X$, is itself again an order ideal: every element of the sum
splits into pieces which remain bounded in absolute value by the
element itself. This yields two decomposition statements, one for
non-negative elements and one with lattice estimates on the absolute
values of the pieces.

\begin{lemma}
\label{lem:orderIdeal-sum}
\lean{OrderIdeal.sum}
\leanok
\uses{def:order-ideal, lem:orderIdeal-ofSolid,
  lem:inf-eq-sub-posPart}
Let $J_1, J_2$ be order ideals of $X$. Their sum $J_1 + J_2$, viewed
as a linear subspace of $X$, is again an order ideal of $X$.
\end{lemma}

\begin{theorem}
\label{thm:orderIdeal-exists-sum-decomp-nonneg}
\lean{OrderIdeal.exists_sum_decomp_nonneg}
\leanok
\uses{lem:orderIdeal-sum, lem:inf-eq-sub-posPart}
Let $J_1, J_2$ be order ideals of $X$ and let $y \in J_1 + J_2$ with
$0 \le y$. Then there exist $y_1, y_2 \in X$ with $y_1 \in J_1$,
$y_2 \in J_2$, $y_1 + y_2 = y$, $0 \le y_1$ and $0 \le y_2$.
\end{theorem}

\begin{theorem}
\label{thm:orderIdeal-exists-sum-decomp}
\lean{OrderIdeal.exists_sum_decomp}
\leanok
\uses{lem:orderIdeal-sum, lem:inf-eq-sub-posPart}
Let $J_1, J_2$ be order ideals of $X$ and let $z \in J_1 + J_2$.
Then there exist $a, b \in X$ with $a \in J_1$, $b \in J_2$,
$a + b = z$, $\lvert a \rvert \le \lvert z \rvert$ and
$\lvert b \rvert \le \lvert z \rvert$.
\end{theorem}

\subsection{The complete lattice of order ideals}

Ordered by inclusion, order ideals of $X$ form a complete lattice:
the trivial subspace $\{0\}$ is the least ideal, the whole space is
the greatest, arbitrary intersections are again ideals, and binary
joins are given by the sum. When $X$ is Archimedean with dimension strictly greater than $1$,
this lattice always contains an ideal distinct from both $\{0\}$ and
$X$.

\begin{lemma}
\label{lem:orderIdeal-complete-lattice}
\lean{OrderIdeal.instTop, OrderIdeal.instBot, OrderIdeal.instInfSet,
  OrderIdeal.instCompleteLattice}
\leanok
\uses{def:order-ideal, lem:orderIdeal-sum,
  lem:orderIdeal-toSubmodule-injective}
The whole space $X$ is the greatest order ideal of $X$, the zero
subspace $\{0\}$ is the least order ideal of $X$, the intersection
of an arbitrary family of order ideals is again an order ideal, and
the order ideals of $X$ ordered by inclusion form a complete lattice
with binary joins given by the sum of ideals.
\end{lemma}

\begin{theorem}
\label{thm:orderIdeal-strongOrderUnit-iff-principal-eq-top}
\lean{OrderIdeal.strongOrderUnit_iff_principal_eq_top}
\leanok
\uses{def:strongOrderUnit, def:orderIdeal-principal,
  lem:orderIdeal-complete-lattice}
Let $X$ be a vector lattice and let $e \in X$. Then $e$ is a strong order
unit of $X$ if and only if $0 \le e$ and the principal ideal $I_e$ equals
$X$.
\end{theorem}

\begin{theorem}
\label{thm:orderIdeal-exists-proper-nontrivial}
\lean{OrderIdeal.exists_proper_nontrivial}
\leanok
\uses{def:orderIdeal-principal, def:is-vl-archimedean,
  thm:exists-pair-ne-zero-isVLDisjoint,
  lem:orderIdeal-complete-lattice}
Let $X$ be an Archimedean vector lattice with dimension strictly
greater than $1$. Then there exists an order ideal $J$ of $X$ with
$J \ne \{0\}$ and $J \ne X$.
\end{theorem}

\subsection{Closure of an ideal}

In a normed vector lattice, the topological closure of an order
ideal is again an order ideal: the solid condition propagates to
the closure by truncating an approximating sequence against the
ambient absolute value. In a Banach lattice this closure operation
gives rise to the bundled notion of closed order ideal, stable
under arbitrary intersections and under binary joins given by the
sum.

\begin{lemma}
\label{lem:orderIdeal-topologicalClosure-solid}
\lean{OrderIdeal.topologicalClosure_solid}
\leanok
\uses{def:order-ideal, def:normed-vector-lattice,
  thm:continuous-sup, thm:continuous-inf, thm:lipschitz-abs}
Let $X$ be a normed vector lattice, let $J$ be an order ideal of
$X$, and let $x, y \in X$ with $x$ in the topological closure of
$J$ and $\lvert y \rvert \le \lvert x \rvert$. Then $y$ lies in the
topological closure of $J$.
\end{lemma}

\begin{lemma}
\label{lem:orderIdeal-topologicalClosure}
\lean{OrderIdeal.topologicalClosure}
\leanok
\uses{lem:orderIdeal-topologicalClosure-solid,
  lem:orderIdeal-ofSolid}
Let $X$ be a normed vector lattice and let $J$ be an order ideal of
$X$. Then the topological closure of $J$ is an order ideal of $X$.
\end{lemma}

\begin{definition}
\label{def:closed-order-ideal}
\lean{ClosedOrderIdeal}
\leanok
\uses{def:order-ideal, def:banach-lattice}
Let $X$ be a Banach lattice. A \emph{closed order ideal} of $X$ is
an order ideal of $X$ which is closed in the norm topology.
\end{definition}

\begin{lemma}
\label{lem:closedOrderIdeal-inf}
\lean{ClosedOrderIdeal.inf}
\leanok
\uses{def:closed-order-ideal}
Let $X$ be a Banach lattice. The intersection of two closed order
ideals of $X$ is a closed order ideal of $X$.
\end{lemma}

\begin{lemma}
\label{lem:closedOrderIdeal-sup}
\lean{ClosedOrderIdeal.sup}
\leanok
\uses{def:closed-order-ideal, lem:orderIdeal-sum,
  thm:orderIdeal-exists-sum-decomp}
Let $X$ be a Banach lattice and let $J_1, J_2$ be closed order
ideals of $X$. Then the sum $J_1 + J_2$ is a closed order ideal of
$X$.
\end{lemma}

\begin{lemma}
\label{lem:closedOrderIdeal-sInf}
\lean{ClosedOrderIdeal.sInf}
\leanok
\uses{def:closed-order-ideal}
Let $X$ be a Banach lattice. The intersection of an arbitrary
family of closed order ideals of $X$ is a closed order ideal of
$X$.
\end{lemma}

\begin{lemma}
\label{lem:closedOrderIdeal-lattice}
\lean{ClosedOrderIdeal.instLattice, ClosedOrderIdeal.instInfSet}
\leanok
\uses{lem:closedOrderIdeal-inf, lem:closedOrderIdeal-sup,
  lem:closedOrderIdeal-sInf}
Let $X$ be a Banach lattice. The closed order ideals of $X$
ordered by inclusion form a lattice with binary joins given by the
sum and binary meets by intersection, and are closed under
arbitrary intersections.
\end{lemma}

\subsection{Quotients}

The quotient of a vector lattice $X$ by an order ideal $J$ carries
a canonical vector lattice structure for which the quotient
projection $X \to X / J$ is a vector lattice homomorphism, and
every vector lattice homomorphism out of $X$ that vanishes on $J$
factors through it. The kernel of a vector lattice homomorphism is
itself an order ideal, so quotients by ideals are precisely the
images of vector lattice quotient maps. Order-theoretic information
lifts well from the quotient: positive elements have positive
representatives, inequalities can be witnessed between chosen
representatives, and increasing positive sequences can be lifted to
increasing positive sequences. The Archimedean property of the
quotient corresponds to a closure condition on $J$ under uniform
convergence. In the normed setting, when $J$ is norm-closed the
quotient norm is a lattice norm, and completeness of $X$ passes to
the quotient: the quotient of a Banach lattice by a closed order
ideal is again a Banach lattice.

\subsubsection{Vector lattice quotient}

Throughout this part $X$ denotes a vector lattice and $J$ an order
ideal of $X$.

\begin{lemma}
\label{lem:orderIdeal-ofKerVecLatHom}
\lean{OrderIdeal.ofKerVecLatHom}
\leanok
\uses{def:order-ideal, lem:orderIdeal-ofSolid,
  lem:abs-eq-zero-iff-zero}
Let $X, Y$ be vector lattices and let $f \colon X \to Y$ be a
vector lattice homomorphism. Then $\ker f$ is an order ideal of
$X$.
\end{lemma}

\begin{lemma}
\label{lem:orderIdeal-instQuotientLattice}
\lean{OrderIdeal.instQuotientLattice}
\leanok
\uses{def:order-ideal}
Let $X$ be a vector lattice and let $J$ be an order ideal of $X$.
The quotient module $X / J$ carries a lattice structure, uniquely
determined by
\[
  [x] \vee [y] = [x \vee y], \qquad [x] \wedge [y] = [x \wedge y],
\]
where $[\cdot]$ denotes the canonical projection $X \to X / J$.
\end{lemma}

\begin{lemma}
\label{lem:orderIdeal-instQuotientIsOrderedAddMonoid}
\lean{OrderIdeal.instQuotientIsOrderedAddMonoid}
\leanok
\uses{lem:orderIdeal-instQuotientLattice}
Let $X$ be a vector lattice and let $J$ be an order ideal of $X$.
The order on the quotient $X / J$ is compatible with addition: for
all $\varphi, \psi, \chi \in X / J$, $\varphi \le \psi$ implies
$\varphi + \chi \le \psi + \chi$.
\end{lemma}

\begin{theorem}
\label{thm:orderIdeal-instQuotientVectorLattice}
\lean{OrderIdeal.instQuotientVectorLattice}
\leanok
\uses{def:vector-lattice, lem:orderIdeal-instQuotientLattice,
  lem:orderIdeal-instQuotientIsOrderedAddMonoid,
  lem:nonneg-smul-sup}
Let $X$ be a vector lattice and let $J$ be an order ideal of $X$.
Then $X / J$ is a vector lattice, with lattice operations,
addition, and scalar multiplication induced from $X$.
\end{theorem}

\begin{lemma}
\label{lem:orderIdeal-mkQVecLatHom}
\lean{OrderIdeal.mkQVecLatHom}
\leanok
\uses{thm:orderIdeal-instQuotientVectorLattice}
Let $X$ be a vector lattice and let $J$ be an order ideal of $X$.
The canonical projection $X \to X / J$ is a vector lattice
homomorphism.
\end{lemma}

\subsubsection{Lifting and factoring through the quotient}

Order-theoretic data on $X / J$ admits canonical lifts to $X$:
positive classes, inequalities, and intervals are all witnessed by
representatives chosen with the corresponding positivity and
ordering. Conversely, linear and lattice homomorphisms out of $X$
that annihilate $J$ factor uniquely through the quotient.

\begin{lemma}
\label{lem:orderIdeal-lift-pos}
\lean{OrderIdeal.lift_pos}
\leanok
\uses{lem:orderIdeal-mkQVecLatHom}
Let $J$ be an order ideal of the vector lattice $X$ and let
$\varphi \in X / J$ with $0 \le \varphi$. Then there exists
$x \in X$ with $0 \le x$ and $[x] = \varphi$.
\end{lemma}

\begin{lemma}
\label{lem:orderIdeal-lift-le}
\lean{OrderIdeal.lift_le}
\leanok
\uses{lem:orderIdeal-mkQVecLatHom}
Let $J$ be an order ideal of the vector lattice $X$, let
$\varphi, \psi \in X / J$ with $\varphi \le \psi$, and let
$x \in X$ with $[x] = \varphi$. Then there exists $y \in X$ with
$[y] = \psi$ and $x \le y$.
\end{lemma}

\begin{lemma}
\label{lem:orderIdeal-lift-between}
\lean{OrderIdeal.lift_between}
\leanok
\uses{lem:orderIdeal-mkQVecLatHom}
Let $J$ be an order ideal of the vector lattice $X$, let
$x, z \in X$ with $x \le z$, and let $\varphi \in X / J$ with
$[x] \le \varphi \le [z]$. Then there exists $y \in X$ with
$[y] = \varphi$, $x \le y$ and $y \le z$.
\end{lemma}

\begin{lemma}
\label{lem:orderIdeal-lift-abs-le}
\lean{OrderIdeal.lift_abs_le}
\leanok
\uses{lem:orderIdeal-lift-between, lem:orderIdeal-mkQVecLatHom}
Let $J$ be an order ideal of the vector lattice $X$, let
$\varphi \in X / J$ and let $y \in X$ with $0 \le y$ and
$[y] = \lvert \varphi \rvert$. Then there exists $x \in X$ with
$[x] = \varphi$ and $\lvert x \rvert \le y$.
\end{lemma}

\begin{lemma}
\label{lem:orderIdeal-lift-increasing-seq}
\lean{OrderIdeal.lift_increasing_seq}
\leanok
\uses{lem:orderIdeal-lift-pos, lem:orderIdeal-lift-le}
Let $J$ be an order ideal of the vector lattice $X$ and let
$(\varphi_n)_{n \in \mathbb{N}}$ be a monotone sequence in $X / J$
with $0 \le \varphi_n$ for every $n$. Then there exists a monotone
sequence $(x_n)_{n \in \mathbb{N}}$ in $X$ with $0 \le x_n$ and
$[x_n] = \varphi_n$ for every $n$.
\end{lemma}

\begin{lemma}
\label{lem:orderIdeal-liftQ-positive}
\lean{OrderIdeal.liftQ_positive}
\leanok
\uses{lem:orderIdeal-lift-pos}
Let $J$ be an order ideal of the vector lattice $X$, let $Y$ be a
vector lattice, and let $T \colon X \to Y$ be a positive linear
operator with $J \subseteq \ker T$. Then the induced linear map
$\bar T \colon X / J \to Y$ is positive.
\end{lemma}

\begin{lemma}
\label{lem:orderIdeal-liftQVecLatHom}
\lean{OrderIdeal.liftQVecLatHom}
\leanok
\uses{lem:orderIdeal-mkQVecLatHom}
Let $J$ be an order ideal of the vector lattice $X$, let $Y$ be a
vector lattice, and let $T \colon X \to Y$ be a vector lattice
homomorphism with $J \subseteq \ker T$. Then $T$ factors through
the quotient as a vector lattice homomorphism
$\bar T \colon X / J \to Y$ with $\bar T \circ [\cdot] = T$.
\end{lemma}

\subsubsection{Uniformly closed ideals and the Archimedean quotient}

The quotient $X / J$ is Archimedean precisely when the ideal $J$ is
closed under uniform convergence: a sequence in $J$ which converges
to $x$ uniformly relative to a fixed positive vector forces $x$ to
lie in $J$.

\begin{definition}
\label{def:orderIdeal-isUniformlyClosed}
\lean{OrderIdeal.IsUniformlyClosed}
\leanok
\uses{def:order-ideal}
Let $J$ be an order ideal of the vector lattice $X$. We say that
$J$ is \emph{uniformly closed} if, for all $x, u \in X$ and all
sequences $(x_n)_{n \in \mathbb{N}}$ in $X$ with $0 \le u$,
$x_n \in J$ for every $n$, and
$\lvert x - x_n \rvert \le \tfrac{1}{n + 1} u$ for every $n$, the
limit $x$ also lies in $J$.
\end{definition}

\begin{theorem}
\label{thm:orderIdeal-quotient-archimedean-iff-uniformlyClosed}
\lean{OrderIdeal.quotient_archimedean_iff_uniformlyClosed}
\leanok
\uses{def:is-vl-archimedean, def:orderIdeal-isUniformlyClosed,
  thm:orderIdeal-instQuotientVectorLattice,
  lem:orderIdeal-lift-pos, lem:abs-eq-zero-iff-zero}
Let $J$ be an order ideal of the vector lattice $X$. Then the
quotient $X / J$ is Archimedean if and only if $J$ is uniformly
closed.
\end{theorem}

\subsubsection{Normed and Banach lattice quotient}

When $X$ is a normed vector lattice and $J$ is a norm-closed order
ideal, the quotient norm is a lattice norm, and so $X / J$ is a
normed vector lattice. If $X$ is a Banach lattice, completeness
passes to the quotient: $X / J$ is a Banach lattice.

\begin{lemma}
\label{lem:orderIdeal-quotientNormedAddCommGroup}
\lean{OrderIdeal.quotientNormedAddCommGroup}
\leanok
\uses{def:normed-vector-lattice, def:order-ideal}
Let $X$ be a normed vector lattice and let $J$ be a norm-closed
order ideal of $X$. Then the quotient norm on $X / J$ is a norm
(i.e., it separates points).
\end{lemma}

\begin{theorem}
\label{thm:orderIdeal-quotientNormedVectorLattice}
\lean{OrderIdeal.quotientNormedVectorLattice}
\leanok
\uses{def:normed-vector-lattice,
  thm:orderIdeal-instQuotientVectorLattice,
  lem:orderIdeal-quotientNormedAddCommGroup,
  lem:orderIdeal-lift-between}
Let $X$ be a normed vector lattice and let $J$ be a norm-closed
order ideal of $X$. Then $X / J$, with the induced lattice and
vector-space structure and the quotient norm, is a normed vector
lattice.
\end{theorem}

\begin{theorem}
\label{thm:orderIdeal-quotientBanachLattice}
\lean{OrderIdeal.quotientBanachLattice}
\leanok
\uses{def:banach-lattice,
  thm:orderIdeal-quotientNormedVectorLattice}
Let $X$ be a Banach lattice and let $J$ be a closed order ideal of
$X$. Then $X / J$, with the induced structures, is a Banach
lattice.
\end{theorem}

\section{Bands}

A \emph{band} of a vector lattice $X$ is an order ideal which is
\emph{order closed}: whenever a non-empty subset of the band has a
supremum in $X$, that supremum lies in the band. Bands form a complete
lattice under inclusion, and to every set $A \subseteq X$ corresponds
the smallest band containing $A$. Two basic constructions appear: the
\emph{disjoint complement} $A^d$ of a set, which is always a band, and
the \emph{band generated} by a set. In an Archimedean vector lattice the
two operations are linked by the identity $A^{dd} = \langle A
\rangle_{\mathrm{band}}$, whose immediate consequences include several
descriptions of the band generated by an order ideal in terms of
suprema, a sequential description of the principal band, the
characterisation of weak order units as those non-negative vectors that
generate the whole space as a band, and norm-closedness of every band
in a normed vector lattice. The next layer is the notion of
\emph{projection band}: a band $B$ such that $X$ decomposes as the
direct sum of $B$ and its disjoint complement. Projection bands form a
Boolean algebra, and the associated band projections are positive
idempotent vector lattice homomorphisms dominated by the identity. A
vector lattice has the \emph{Projection Property} (PP) when every band
is a projection band, and the \emph{Principal Projection Property}
(PPP) when every principal band is a projection band; (sigma) order
completeness implies (PPP) PP, and either implies the Archimedean
property. The section closes with the structural theorem that, under
PPP, the principal-band projections along a maximal disjoint family
exhibit $X$ as an order-dense vector lattice subobject of the product
of the principal bands.

\subsection{Bands}

Throughout $X$ denotes a vector lattice.

\begin{definition}
\label{def:band}
\lean{Band}
\leanok
\uses{def:order-ideal}
Let $X$ be a vector lattice. A \emph{band} of $X$ is an order ideal
$B$ of $X$ which is \emph{order closed}: for every non-empty subset
$S \subseteq B$ and every $x \in X$, if $x$ is the supremum of $S$ in
$X$, then $x \in B$.
\end{definition}

\begin{lemma}
\label{lem:band-sSup-mem}
\lean{Band.sSup_mem}
\leanok
\uses{def:band}
Let $B$ be a band of $X$, let $S \subseteq B$ be non-empty, and let
$x \in X$ with $x = \sup S$ in $X$. Then $x \in B$.
\end{lemma}

\begin{lemma}
\label{lem:band-ofPosDirectedSSupMem}
\lean{Band.ofPosDirectedSSupMem}
\leanok
\uses{def:band}
Let $J$ be an order ideal of $X$ such that for every non-empty
upward-directed subset $S \subseteq J$ of non-negative elements with a
supremum $x$ in $X$, $x \in J$. Then $J$ is a band.
\end{lemma}

\begin{lemma}
\label{lem:band-instConditionallyCompleteLatticeSubtype}
\lean{Band.instConditionallyCompleteLatticeSubtype}
\leanok
\uses{def:band, thm:conditionallyCompleteLatticeOfPosSet}
Let $X$ be a sigma Dedekind complete vector lattice and let $B$ be a
band of $X$. Then $B$ is itself a sigma Dedekind complete vector
lattice under the order inherited from $X$.
\end{lemma}

\subsection{The complete lattice of bands}

Ordered by inclusion, bands of $X$ form a complete lattice: the whole
space is the greatest band, the trivial subspace $\{0\}$ is the least,
arbitrary intersections of bands are bands, and binary meets are given
by intersection.

\begin{lemma}
\label{lem:band-complete-lattice}
\lean{Band.instTop, Band.instBot, Band.instInfSet,
  Band.instCompleteLattice}
\leanok
\uses{def:band}
The whole space $X$ is the greatest band of $X$, the zero subspace
$\{0\}$ is the least band of $X$, the intersection of an arbitrary
family of bands is again a band, and the bands of $X$ ordered by
inclusion form a complete lattice with binary meets given by
intersection.
\end{lemma}

\subsection{Disjoint complements}

The \emph{disjoint complement} $A^d$ of a set $A \subseteq X$ is the
set of elements of $X$ that are disjoint from every member of $A$. It
is always a band, behaves antimonotonically under inclusion, satisfies
the triple-complement identity $A^{ddd} = A^d$, and converts unions
into intersections. In a normed vector lattice it is moreover norm
closed.

\begin{definition}
\label{def:disjointComplement}
\lean{disjointComplement}
\leanok
\uses{def:isVLDisjoint}
Let $X$ be a vector lattice and let $A \subseteq X$. The \emph{disjoint
complement} of $A$ is
\[
  A^d := \{x \in X : \text{$x$ is disjoint from every $a \in A$}\}.
\]
\end{definition}

\begin{lemma}
\label{lem:disjointComplement-anti}
\lean{disjointComplement_anti}
\leanok
\uses{def:disjointComplement}
Let $A, B \subseteq X$ with $A \subseteq B$. Then $B^d \subseteq A^d$.
\end{lemma}

\begin{lemma}
\label{lem:disjointComplement-inter-eq-zero}
\lean{disjointComplement_inter_eq_zero}
\leanok
\uses{def:disjointComplement, lem:abs-eq-zero-iff-zero}
Let $A \subseteq X$. Then $A \cap A^d \subseteq \{0\}$.
\end{lemma}

\begin{lemma}
\label{lem:subset-disjointComplement-disjointComplement}
\lean{subset_disjointComplement_disjointComplement}
\leanok
\uses{def:disjointComplement}
Every $A \subseteq X$ satisfies $A \subseteq A^{dd}$.
\end{lemma}

\begin{lemma}
\label{lem:disjointComplement-disjointComplement-disjointComplement}
\lean{disjointComplement_disjointComplement_disjointComplement}
\leanok
\uses{lem:subset-disjointComplement-disjointComplement,
  lem:disjointComplement-anti}
Every $A \subseteq X$ satisfies $A^{ddd} = A^d$.
\end{lemma}

\begin{lemma}
\label{lem:disjointComplement-union}
\lean{disjointComplement_union}
\leanok
\uses{def:disjointComplement}
For all $A, B \subseteq X$, $(A \cup B)^d = A^d \cap B^d$.
\end{lemma}

\begin{lemma}
\label{lem:band-disjointComplement}
\lean{Band.disjointComplement}
\leanok
\uses{def:band, def:disjointComplement, lem:band-ofPosDirectedSSupMem}
Let $A \subseteq X$. Then $A^d$ is a band of $X$.
\end{lemma}

\begin{lemma}
\label{lem:isClosed-disjointComplement}
\lean{isClosed_disjointComplement}
\leanok
\uses{def:disjointComplement, def:normed-vector-lattice,
  thm:lipschitz-abs, thm:continuous-inf}
Let $X$ be a normed vector lattice and let $A \subseteq X$. Then $A^d$
is norm closed in $X$.
\end{lemma}

\subsection{The band generated by a set}

For any subset $A \subseteq X$ there is a smallest band of $X$
containing $A$, the \emph{band generated} by $A$. In an Archimedean
vector lattice it coincides with the double disjoint complement
$A^{dd}$. From this one obtains four equivalent descriptions of the
band generated by an order ideal in terms of suprema, a concrete
sequential description of the principal band, the characterisation of
weak order units as non-negative generators of the whole space, and
norm-closedness of every band in a normed vector lattice.

\begin{definition}
\label{def:band-generated}
\lean{Band.generated}
\leanok
\uses{def:band, lem:band-complete-lattice}
Let $A \subseteq X$. The \emph{band generated} by $A$ is the
intersection of all bands of $X$ containing $A$.
\end{definition}

\begin{lemma}
\label{lem:band-subset-generated}
\lean{Band.subset_generated}
\leanok
\uses{def:band-generated}
For every $A \subseteq X$, $A$ is contained in the band generated by
$A$.
\end{lemma}

\begin{lemma}
\label{lem:band-generated-le}
\lean{Band.generated_le}
\leanok
\uses{def:band-generated}
Let $A \subseteq X$ and let $B$ be a band of $X$ with $A \subseteq B$.
Then the band generated by $A$ is contained in $B$.
\end{lemma}

\begin{theorem}
\label{thm:band-disjointComplement-disjointComplement-eq-generated}
\lean{Band.disjointComplement_disjointComplement_eq_generated}
\leanok
\uses{def:band-generated, def:disjointComplement,
  def:is-vl-archimedean, lem:band-disjointComplement,
  lem:disjointComplement-inter-eq-zero}
Let $X$ be an Archimedean vector lattice and let $A \subseteq X$. Then
the band generated by $A$ equals $A^{dd}$.
\end{theorem}

\begin{theorem}
\label{thm:band-eq-disjointComplement-disjointComplement}
\lean{Band.eq_disjointComplement_disjointComplement}
\leanok
\uses{thm:band-disjointComplement-disjointComplement-eq-generated}
Let $X$ be an Archimedean vector lattice and let $B$ be a band of $X$.
Then $B = B^{dd}$.
\end{theorem}

\begin{theorem}
\label{thm:band-exists-eq-disjointComplement}
\lean{Band.exists_eq_disjointComplement}
\leanok
\uses{thm:band-eq-disjointComplement-disjointComplement}
Let $X$ be an Archimedean vector lattice. Every band of $X$ is the
disjoint complement of some subset of $X$.
\end{theorem}

\begin{theorem}
\label{thm:band-mem-generated-orderIdeal-iff-isLUB-directed}
\lean{Band.mem_generated_orderIdeal_iff_isLUB_directed}
\leanok
\uses{def:band-generated, def:order-ideal, def:is-vl-archimedean,
  thm:band-disjointComplement-disjointComplement-eq-generated}
Let $X$ be an Archimedean vector lattice, let $J$ be an order ideal of
$X$, and let $x \in X$ with $0 \le x$. Then $x$ belongs to the band
generated by $J$ if and only if there exists a non-empty
upward-directed subset $S \subseteq J$ of non-negative elements with
$x = \sup S$ in $X$.
\end{theorem}

\begin{theorem}
\label{thm:band-mem-generated-orderIdeal-iff-isLUB-nonneg}
\lean{Band.mem_generated_orderIdeal_iff_isLUB_nonneg}
\leanok
\uses{def:band-generated, def:order-ideal, def:is-vl-archimedean,
  thm:band-disjointComplement-disjointComplement-eq-generated}
Let $X$ be an Archimedean vector lattice, let $J$ be an order ideal of
$X$, and let $x \in X$ with $0 \le x$. Then $x$ belongs to the band
generated by $J$ if and only if there exists a non-empty subset
$D \subseteq J$ of non-negative elements with $x = \sup D$ in $X$.
\end{theorem}

\begin{theorem}
\label{thm:band-mem-generated-orderIdeal-iff-isLUB-inter}
\lean{Band.mem_generated_orderIdeal_iff_isLUB_inter}
\leanok
\uses{def:band-generated, def:order-ideal, def:is-vl-archimedean,
  thm:band-disjointComplement-disjointComplement-eq-generated}
Let $X$ be an Archimedean vector lattice, let $J$ be an order ideal of
$X$, and let $x \in X$ with $0 \le x$. Then $x$ belongs to the band
generated by $J$ if and only if $x = \sup (J \cap [0, x])$ in $X$.
\end{theorem}

\begin{theorem}
\label{thm:band-mem-generated-singleton-iff-isLUB-inf-nsmul}
\lean{Band.mem_generated_singleton_iff_isLUB_inf_nsmul}
\leanok
\uses{def:band-generated, def:is-vl-archimedean,
  thm:band-disjointComplement-disjointComplement-eq-generated}
Let $X$ be an Archimedean vector lattice and let $a, x \in X$ with
$0 \le a$ and $0 \le x$. Then $x$ belongs to the band generated by
$\{a\}$ if and only if
\[
  x = \sup_{n \in \mathbb{N}} (x \wedge n a).
\]
\end{theorem}

\begin{theorem}
\label{thm:weakOrderUnit-iff-generated-singleton-eq-top}
\lean{weakOrderUnit_iff_generated_singleton_eq_top}
\leanok
\uses{def:weakOrderUnit, def:band-generated, def:is-vl-archimedean,
  thm:band-disjointComplement-disjointComplement-eq-generated}
Let $X$ be an Archimedean vector lattice and let $e \in X$ with
$0 \le e$. Then $e$ is a weak order unit of $X$ if and only if the
band generated by $\{e\}$ is the whole space $X$.
\end{theorem}

\begin{theorem}
\label{thm:band-isClosed-coe}
\lean{Band.isClosed_coe}
\leanok
\uses{def:band, def:normed-vector-lattice,
  thm:band-eq-disjointComplement-disjointComplement,
  lem:isClosed-disjointComplement}
Let $X$ be a normed vector lattice. Every band of $X$ is norm closed.
\end{theorem}

\begin{lemma}
\label{lem:band-instBanachLatticeSubtype}
\lean{Band.instBanachLatticeSubtype}
\leanok
\uses{def:band, def:banach-lattice, thm:band-isClosed-coe,
  lem:vectorSublattice-instBanachLatticeSubtype}
Let $X$ be a Banach lattice and let $B$ be a band of $X$. Then $B$,
equipped with the inherited norm and lattice operations, is a Banach
lattice.
\end{lemma}

\subsection{Projection bands}

A band $B$ of $X$ is a \emph{projection band} when every $x \in X$
admits a (necessarily unique) decomposition $x = y + z$ with $y \in B$
and $z \in B^d$. The associated \emph{band projection}
$P_B \colon X \to X$ sending $x$ to its $B$-component is a positive
idempotent vector lattice homomorphism dominated by the identity, and
every linear operator $T$ on $X$ with $T^2 = T$ and $0 \le T \le I$
arises in this way. Projection bands of $X$ form a Boolean algebra
under inclusion, with complement given by the disjoint complement and
binary meets and joins compatible with the corresponding band
projections.

\begin{definition}
\label{def:projectionBand}
\lean{ProjectionBand}
\leanok
\uses{def:band, def:disjointComplement}
Let $X$ be a vector lattice. A \emph{projection band} of $X$ is a
band $B$ of $X$ such that for every $x \in X$ there exist $y, z \in X$
with $y \in B$, $z \in B^d$ and $x = y + z$.
\end{definition}

\begin{lemma}
\label{lem:projectionBand-decomposition}
\lean{ProjectionBand.decomposition}
\leanok
\uses{def:projectionBand}
Let $B$ be a projection band of $X$. For every $x \in X$ there exist
$y, z \in X$ with $y \in B$, $z \in B^d$ and $x = y + z$.
\end{lemma}

\begin{lemma}
\label{lem:projectionBand-decomposition-unique}
\lean{ProjectionBand.decomposition_unique}
\leanok
\uses{def:projectionBand, lem:disjointComplement-inter-eq-zero}
Let $B$ be a projection band of $X$ and let $x \in X$. If $x = y_1 +
z_1 = y_2 + z_2$ with $y_1, y_2 \in B$ and $z_1, z_2 \in B^d$, then
$y_1 = y_2$ and $z_1 = z_2$.
\end{lemma}

\begin{definition}
\label{def:projectionBand-bandProjection}
\lean{ProjectionBand.bandProjection}
\leanok
\uses{def:projectionBand, lem:projectionBand-decomposition,
  lem:projectionBand-decomposition-unique}
Let $B$ be a projection band of $X$. The \emph{band projection}
associated to $B$ is the linear map $P_B \colon X \to X$ sending
$x$ to its $B$-component in the decomposition $x = y + z$ with
$y \in B$ and $z \in B^d$.
\end{definition}

\begin{lemma}
\label{lem:projectionBand-bandProjection-sq}
\lean{ProjectionBand.bandProjection_sq}
\leanok
\uses{def:projectionBand-bandProjection}
Let $B$ be a projection band of $X$. Then $P_B \circ P_B = P_B$.
\end{lemma}

\begin{lemma}
\label{lem:projectionBand-bandProjection-mem}
\lean{ProjectionBand.bandProjection_mem}
\leanok
\uses{def:projectionBand-bandProjection}
Let $B$ be a projection band of $X$. For every $x \in X$, $P_B(x) \in
B$.
\end{lemma}

\begin{lemma}
\label{lem:projectionBand-id-sub-bandProjection-mem}
\lean{ProjectionBand.id_sub_bandProjection_mem}
\leanok
\uses{def:projectionBand-bandProjection}
Let $B$ be a projection band of $X$. For every $x \in X$, $x - P_B(x)
\in B^d$.
\end{lemma}

\begin{lemma}
\label{lem:projectionBand-bandProjection-nonneg}
\lean{ProjectionBand.bandProjection_nonneg}
\leanok
\uses{def:projectionBand-bandProjection}
Let $B$ be a projection band of $X$. Then $P_B$ is a positive
operator: $0 \le P_B(x)$ for every $0 \le x$.
\end{lemma}

\begin{lemma}
\label{lem:projectionBand-bandProjection-le-id}
\lean{ProjectionBand.bandProjection_le_id}
\leanok
\uses{def:projectionBand-bandProjection,
  lem:projectionBand-bandProjection-nonneg}
Let $B$ be a projection band of $X$. Then $P_B \le \mathrm{id}_X$
pointwise on $X_+$: $P_B(x) \le x$ for every $0 \le x$.
\end{lemma}

\begin{lemma}
\label{lem:projectionBand-range-bandProjection}
\lean{ProjectionBand.range_bandProjection}
\leanok
\uses{def:projectionBand-bandProjection,
  lem:projectionBand-bandProjection-mem}
Let $B$ be a projection band of $X$. Then the range of $P_B$ equals
$B$.
\end{lemma}

\begin{lemma}
\label{lem:projectionBand-disjointComplementProjectionBand}
\lean{ProjectionBand.disjointComplementProjectionBand}
\leanok
\uses{def:projectionBand, lem:band-disjointComplement,
  lem:projectionBand-decomposition}
Let $B$ be a projection band of $X$. Then $B^d$ is also a projection
band of $X$.
\end{lemma}

\begin{theorem}
\label{thm:projectionBand-iff-add-disjointComplement}
\lean{ProjectionBand.projectionBand_iff_add_disjointComplement}
\leanok
\uses{def:projectionBand, def:order-ideal,
  lem:disjointComplement-inter-eq-zero}
Let $J$ be an order ideal of $X$. Then $J$ is a projection band of
$X$ if and only if every $x \in X$ can be written as $x = y + z$ with
$y \in J$ and $z \in J^d$.
\end{theorem}

\begin{theorem}
\label{thm:projectionBand-of-direct-sum-of-disjoint}
\lean{ProjectionBand.of_direct_sum_of_disjoint}
\leanok
\uses{def:projectionBand, def:order-ideal, def:disjointComplement,
  thm:projectionBand-iff-add-disjointComplement}
Let $J_1, J_2$ be order ideals of $X$ such that every element of
$J_1$ is disjoint from every element of $J_2$ and every $x \in X$
admits a decomposition $x = y + z$ with $y \in J_1$ and $z \in J_2$.
Then $J_1^d = J_2$ and $J_1$ is a projection band of $X$.
\end{theorem}

\begin{theorem}
\label{thm:projectionBand-bandProjection-iff}
\lean{ProjectionBand.bandProjection_iff}
\leanok
\uses{def:projectionBand-bandProjection,
  lem:projectionBand-bandProjection-sq,
  lem:projectionBand-bandProjection-nonneg,
  lem:projectionBand-bandProjection-le-id}
Let $T \colon X \to X$ be a linear operator. Then $T$ is the band
projection of some projection band of $X$ if and only if $T \circ T =
T$, $0 \le T$ and $T \le \mathrm{id}_X$ (in the pointwise order on
positive operators).
\end{theorem}

\begin{lemma}
\label{lem:projectionBand-bandProjection-isVecLatHom}
\lean{ProjectionBand.bandProjection_isVecLatHom}
\leanok
\uses{def:projectionBand-bandProjection,
  lem:projectionBand-bandProjection-sq,
  lem:projectionBand-bandProjection-nonneg,
  lem:projectionBand-bandProjection-le-id}
Let $B$ be a projection band of $X$. Then $P_B$ is a vector lattice
homomorphism: it preserves $\vee$ and $\wedge$.
\end{lemma}

\begin{lemma}
\label{lem:projectionBand-booleanAlgebra}
\lean{ProjectionBand.instTop, ProjectionBand.instBot,
  ProjectionBand.instCompl, ProjectionBand.instMin,
  ProjectionBand.instMax, ProjectionBand.instBooleanAlgebra}
\leanok
\uses{def:projectionBand, lem:projectionBand-disjointComplementProjectionBand}
The whole space $X$ and the trivial subspace $\{0\}$ are projection
bands of $X$. The disjoint complement of a projection band is a
projection band, and the meet $B_1 \cap B_2$ and the
direct-sum-style join $\{y + z : y \in B_1, z \in B_2\}$ of two
projection bands are projection bands. With these operations the
projection bands of $X$ form a Boolean algebra.
\end{lemma}

\begin{lemma}
\label{lem:projectionBand-coe-compl}
\lean{ProjectionBand.coe_compl}
\leanok
\uses{lem:projectionBand-booleanAlgebra, def:disjointComplement}
For every projection band $B$ of $X$, the complement $B^c$ in the
Boolean algebra of projection bands equals $B^d$.
\end{lemma}

\begin{lemma}
\label{lem:projectionBand-coe-band-sup-eq-add}
\lean{ProjectionBand.coe_band_sup_eq_add}
\leanok
\uses{lem:projectionBand-booleanAlgebra, lem:band-complete-lattice}
For all projection bands $B_1, B_2$ of $X$, the band $B_1 \vee B_2$
equals $\{y + z : y \in B_1,\, z \in B_2\}$.
\end{lemma}

\begin{lemma}
\label{lem:projectionBand-bandProjection-inf}
\lean{ProjectionBand.bandProjection_inf}
\leanok
\uses{lem:projectionBand-booleanAlgebra,
  def:projectionBand-bandProjection}
For all projection bands $B_1, B_2$ of $X$,
$P_{B_1 \cap B_2} = P_{B_1} \circ P_{B_2}$.
\end{lemma}

\begin{lemma}
\label{lem:projectionBand-bandProjection-comm}
\lean{ProjectionBand.bandProjection_comm}
\leanok
\uses{lem:projectionBand-bandProjection-inf}
For all projection bands $B_1, B_2$ of $X$,
$P_{B_1} \circ P_{B_2} = P_{B_2} \circ P_{B_1}$.
\end{lemma}

\begin{lemma}
\label{lem:projectionBand-bandProjection-sup}
\lean{ProjectionBand.bandProjection_sup}
\leanok
\uses{lem:projectionBand-booleanAlgebra,
  def:projectionBand-bandProjection,
  lem:projectionBand-bandProjection-inf}
For all projection bands $B_1, B_2$ of $X$,
\[
  P_{B_1 \vee B_2} = P_{B_1} + P_{B_2} - P_{B_1} \circ P_{B_2}.
\]
\end{lemma}

\begin{lemma}
\label{lem:projectionBand-bandProjection-compl}
\lean{ProjectionBand.bandProjection_compl}
\leanok
\uses{lem:projectionBand-booleanAlgebra,
  def:projectionBand-bandProjection}
For every projection band $B$ of $X$,
$P_{B^c} = \mathrm{id}_X - P_B$.
\end{lemma}

\begin{theorem}
\label{thm:projectionBand-isLUB-inter-Icc-bandProjection}
\lean{ProjectionBand.isLUB_inter_Icc_bandProjection}
\leanok
\uses{def:projectionBand-bandProjection, def:is-vl-archimedean}
Let $X$ be an Archimedean vector lattice, let $B$ be a projection
band of $X$, and let $x \in X$ with $0 \le x$. Then $P_B(x)$ is the
supremum of $B \cap [0, x]$ in $X$.
\end{theorem}

\begin{theorem}
\label{thm:projectionBand-projectionBand-iff-isLUB-inter-Icc}
\lean{ProjectionBand.projectionBand_iff_isLUB_inter_Icc}
\leanok
\uses{def:band, def:projectionBand, def:is-vl-archimedean,
  thm:projectionBand-isLUB-inter-Icc-bandProjection}
Let $X$ be an Archimedean vector lattice and let $B$ be a band of
$X$. Then $B$ is a projection band of $X$ if and only if for every
$x \in X$ with $0 \le x$, the set $B \cap [0, x]$ admits a supremum
in $X$.
\end{theorem}

\begin{theorem}
\label{thm:projectionBand-isLUB-range-bandProjection-of-principalBand}
\lean{ProjectionBand.isLUB_range_bandProjection_of_principalBand}
\leanok
\uses{def:projectionBand-bandProjection, def:band-generated,
  def:is-vl-archimedean,
  thm:band-mem-generated-singleton-iff-isLUB-inf-nsmul}
Let $X$ be an Archimedean vector lattice, let $a \in X$ with
$0 \le a$, and let $B$ be a projection band of $X$ that equals the
band generated by $\{a\}$. Then for every $x \in X$ with $0 \le x$,
\[
  P_B(x) = \sup_{n \in \mathbb{N}} (x \wedge n a).
\]
\end{theorem}

\begin{theorem}
\label{thm:projectionBand-principalBand-projectionBand-iff-isLUB-range}
\lean{ProjectionBand.principalBand_projectionBand_iff_isLUB_range}
\leanok
\uses{def:band-generated, def:projectionBand, def:is-vl-archimedean,
  thm:projectionBand-projectionBand-iff-isLUB-inter-Icc,
  thm:band-mem-generated-singleton-iff-isLUB-inf-nsmul}
Let $X$ be an Archimedean vector lattice and let $a \in X$ with
$0 \le a$. Then the band generated by $\{a\}$ is a projection band of
$X$ if and only if for every $x \in X$ with $0 \le x$, the sequence
$(x \wedge n a)_{n \in \mathbb{N}}$ admits a supremum in $X$.
\end{theorem}

\subsection{The projection properties}

A vector lattice has the \emph{Projection Property} (PP) when every
band is a projection band, and the \emph{Principal Projection
Property} (PPP) when every principal band — that is, every band
generated by a single element — is a projection band. Order
completeness implies PP and sigma order completeness implies PPP. Both
properties imply the Archimedean property.

\begin{definition}
\label{def:hasProjectionProperty}
\lean{HasProjectionProperty}
\leanok
\uses{def:band, def:projectionBand}
A vector lattice $X$ has the \emph{Projection Property} (PP) when
every band of $X$ is a projection band.
\end{definition}

\begin{definition}
\label{def:hasPrincipalProjectionProperty}
\lean{HasPrincipalProjectionProperty}
\leanok
\uses{def:band-generated, def:projectionBand}
A vector lattice $X$ has the \emph{Principal Projection Property}
(PPP) when for every $a \in X$, the band generated by $\{a\}$ is a
projection band of $X$.
\end{definition}

\begin{lemma}
\label{lem:hasPrincipalProjectionProperty-of-hasProjectionProperty}
\lean{HasPrincipalProjectionProperty.of_hasProjectionProperty}
\leanok
\uses{def:hasProjectionProperty, def:hasPrincipalProjectionProperty}
Every vector lattice with the Projection Property has the Principal
Projection Property.
\end{lemma}

\begin{theorem}
\label{thm:hasProjectionProperty-of-isOrderComplete}
\lean{HasProjectionProperty.of_isOrderComplete}
\leanok
\uses{def:hasProjectionProperty,
  thm:projectionBand-projectionBand-iff-isLUB-inter-Icc}
Every Dedekind complete vector lattice has the Projection Property.
\end{theorem}

\begin{theorem}
\label{thm:hasPrincipalProjectionProperty-of-isSigmaOrderComplete}
\lean{HasPrincipalProjectionProperty.of_isSigmaOrderComplete}
\leanok
\uses{def:hasPrincipalProjectionProperty,
  thm:projectionBand-principalBand-projectionBand-iff-isLUB-range}
Every sigma Dedekind complete vector lattice has the Principal
Projection Property.
\end{theorem}

\begin{theorem}
\label{thm:exists-disjoint-sum-eq-sup-of-hasPrincipalProjectionProperty}
\lean{exists_disjoint_sum_eq_sup_of_hasPrincipalProjectionProperty}
\uses{def:hasPrincipalProjectionProperty, def:isVLDisjoint}
Let $X$ be a vector lattice with the Principal Projection Property,
let $n \in \mathbb{N}$, and let $x_0, \dots, x_n \in X$ with
$0 \le x_i$ for every $i$. There exist $y_0, \dots, y_n \in X$ which
are pairwise disjoint, satisfy $0 \le y_i \le x_i$ for every $i$, and
fulfil
\[
  \sum_{i=0}^{n} y_i = \bigvee_{i=0}^{n} x_i.
\]
\end{theorem}

\begin{theorem}
\label{thm:isVLArchimedean-of-hasPrincipalProjectionProperty}
\lean{isVLArchimedean_of_hasPrincipalProjectionProperty}
\uses{def:hasPrincipalProjectionProperty, def:is-vl-archimedean}
Every vector lattice with the Principal Projection Property is
Archimedean.
\end{theorem}

\begin{theorem}
\label{thm:isVLArchimedean-of-hasProjectionProperty}
\lean{isVLArchimedean_of_hasProjectionProperty}
\leanok
\uses{thm:isVLArchimedean-of-hasPrincipalProjectionProperty,
  lem:hasPrincipalProjectionProperty-of-hasProjectionProperty}
Every vector lattice with the Projection Property is Archimedean.
\end{theorem}

\subsection{Decomposition along a maximal disjoint family}

Under the Principal Projection Property, every principal band carries a
canonical band projection. Choosing a maximal disjoint family $\Lambda
\subseteq X$ of strictly positive elements, every positive vector of
$X$ is the supremum of its principal-band projections along $\Lambda$,
and the map sending $x$ to the tuple of its principal-band projections
realises $X$ as an order-dense vector lattice subobject of the product
of the principal bands.

\begin{definition}
\label{def:band-principalProjectionBand}
\lean{Band.principalProjectionBand}
\leanok
\uses{def:hasPrincipalProjectionProperty, def:projectionBand}
Let $X$ be a vector lattice with the Principal Projection Property
and let $a \in X$. We denote by $B_a$ a chosen projection band of $X$
equal to the band generated by $\{a\}$.
\end{definition}

\begin{lemma}
\label{lem:band-principalProjectionBand-coe}
\lean{Band.principalProjectionBand_coe}
\leanok
\uses{def:band-principalProjectionBand, def:band-generated}
Let $X$ be a vector lattice with the Principal Projection Property
and let $a \in X$. Then $B_a$ equals the band generated by $\{a\}$.
\end{lemma}

\begin{definition}
\label{def:band-principalBandProjection}
\lean{Band.principalBandProjection}
\leanok
\uses{def:band-principalProjectionBand,
  def:projectionBand-bandProjection}
Let $X$ be a vector lattice with the Principal Projection Property
and let $a \in X$. The \emph{principal band projection} associated to
$a$ is the band projection $P_a := P_{B_a} \colon X \to X$ associated
to the principal projection band $B_a$.
\end{definition}

\begin{theorem}
\label{thm:isLUB-principalBandProjection-of-isMaximalDisjoint}
\lean{isLUB_principalBandProjection_of_isMaximalDisjoint}
\leanok
\uses{def:band-principalBandProjection, def:isMaximalDisjoint,
  thm:isMaximalDisjoint-iff-forall-eq-zero,
  thm:isVLArchimedean-of-hasPrincipalProjectionProperty,
  lem:projectionBand-id-sub-bandProjection-mem,
  lem:projectionBand-bandProjection-le-id}
Let $X$ be a vector lattice with the Principal Projection Property
and let $\Lambda \subseteq X$ be a maximal disjoint family with
$0 < a$ for every $a \in \Lambda$. For every $x \in X$ with
$0 \le x$,
\[
  x = \sup_{a \in \Lambda} P_a(x).
\]
\end{theorem}

\begin{theorem}
\label{thm:exists-orderDense-lattice-embedding-of-isMaximalDisjoint}
\lean{exists_orderDense_lattice_embedding_of_isMaximalDisjoint}
\leanok
\uses{def:band-principalBandProjection, def:isMaximalDisjoint,
  thm:isLUB-principalBandProjection-of-isMaximalDisjoint,
  thm:isVLArchimedean-of-hasPrincipalProjectionProperty,
  lem:projectionBand-bandProjection-isVecLatHom}
Let $X$ be a vector lattice with the Principal Projection Property
and let $\Lambda \subseteq X$ be a maximal disjoint family with
$0 < a$ for every $a \in \Lambda$. The map
\[
  T \colon X \to \prod_{a \in \Lambda} B_a, \qquad
  T(x) := (P_a(x))_{a \in \Lambda},
\]
is an injective vector lattice homomorphism whose range is order
dense in the product of the principal bands.
\end{theorem}

